%% cv_altacv.tex — Dipayan Sarkar (AltaCV style)
\documentclass[10pt, a4paper, ragged2e, withhyper]{altacv}

\geometry{left=1.2cm, right=1.2cm, top=1.5cm, bottom=1.5cm, columnsep=1cm}

\usepackage{paracol}
\usepackage{hyperref}

% Colors
\definecolor{SlateGrey}{HTML}{2E2E2E}
\definecolor{LightGrey}{HTML}{666666}
\definecolor{DarkPastelBlue}{HTML}{2B6CB0}
\definecolor{PastelBlue}{HTML}{3182CE}
\colorlet{name}{black}
\colorlet{tagline}{DarkPastelBlue}
\colorlet{heading}{DarkPastelBlue}
\colorlet{headingrule}{DarkPastelBlue}
\colorlet{subheading}{PastelBlue}
\colorlet{accent}{DarkPastelBlue}
\colorlet{emphasis}{SlateGrey}
\colorlet{body}{SlateGrey}

\renewcommand{\namefont}{\Huge\bfseries}
\renewcommand{\personalinfofont}{\small}
\renewcommand{\cvsectionfont}{\large\bfseries}
\NewInfoField{scholar}{\faGraduationCap}[https://scholar.google.com/]

\begin{document}

% ---------- Header ----------
\name{Dipayan Sarkar}{}
\tagline{PhD Researcher in Bioinformatics \& Computational Biology}
\photoR{2.5cm}{dpn_photo}

\personalinfo{%
  \email{dipayansarkar26@gmail.com}
  \phone{+91 9735490241}
  \location{Siliguri, West Bengal, India}
  \homepage{dipayansarkar.com}
  \github{Dipayan26}
  \linkedin{dipayansarkar}
  \scholar{Google Scholar}
}

\makecvheader

\columnratio{0.6}

\begin{paracol}{2}

% ===== LEFT COLUMN =====

\cvsection{Research Summary}

Currently pursuing a PhD in bioinformatics and computational biology, focused on deep learning methods for biological sequence-based prediction and generative model development. Experience includes neural network training, transfer learning with pretrained protein language models, and analysis of large-scale biological datasets. Current work involves autoregressive generative modeling of biological sequences, with ongoing exploration of diffusion-based approaches.

\medskip

\cvsection{Work Experience}

\cvevent{PhD Research Scholar}{University of North Bengal, Dept. of Bioinformatics}{2022 -- Present}{Siliguri, India}
\begin{itemize}
  \item Designed and implemented deep learning architectures for sequence-based protein--protein interaction prediction.
  \item Developed attention-based hybrid models integrating transfer learning from pretrained protein language models.
  \item Built and deployed web platforms for model inference and protein interaction network visualization.
\end{itemize}

\cvsection{Publications}

\begin{enumerate}
  \item Sarkar, D., \& Sarkar, C. (2025). AttnSeq-PPI: Enhancing protein-protein interaction network prediction using transfer learning-driven hybrid attention. \textit{BBA---Proteins and Proteomics}, 141102.
  \item Sarkar, C.; Sarkar, D.; Parsad, R.; Mishra, D. (2022). Package \texttt{EGRNi}. CRAN.
\end{enumerate}

\cvsection{Software \& Tools}

\cvachievement{\faCode}{AttnSeq-PPI}{Deep learning-based PPI prediction platform --- \href{https://compbiosysnbu.in/attnseqppi/}{\texttt{compbiosysnbu.in/attnseqppi/}}}

\cvachievement{\faCode}{ARACoFusion-PPI}{Arabidopsis-specific PPI prediction tool --- \href{https://aracofusion.compbiosysnbu.in/}{\texttt{aracofusion.compbiosysnbu.in/}}}

% ===== RIGHT COLUMN =====
\switchcolumn

\cvsection{Education}

\cvevent{Ph.D. Bioinformatics}{University of North Bengal}{Expected Mar 2027}{}
\begin{itemize}
  \item Advisor: Dr. Chiranjib Sarkar
  \item Focus: Deep learning for PPI network prediction
\end{itemize}

\divider

\cvevent{M.Sc. Botany}{University of North Bengal}{2021}{}
Specialization: Biochemistry \\ 77.25\% (First Class)

\divider

\cvevent{B.Sc. (Hons.) Botany}{Ananda Chandra College, UNB}{2019}{}
63.50\% (First Class)

\divider

\cvevent{Higher Secondary (XII)}{Haldibari High School, WBCHSE}{2016}{}
80.00\% (First Class)

\divider

\cvevent{Secondary (X)}{Haldibari High School, WBBSE}{2014}{}
91.58\% (First Class)

\cvsection{Skills}

\cvtag{Python} \cvtag{PyTorch} \cvtag{HuggingFace}
\cvtag{R} \cvtag{C}

\medskip
\cvtag{Transformers} \cvtag{Attention Mechanisms}
\cvtag{Transfer Learning} \cvtag{Sequence Modeling}

\medskip
\cvtag{Autoregressive Models} \cvtag{Diffusion Models}

\medskip
\cvtag{GPU Training} \cvtag{HPC} \cvtag{Bioinformatics}

\medskip
\cvtag{Docker} \cvtag{Flask} \cvtag{FastAPI}
\cvtag{NumPy} \cvtag{Pandas} \cvtag{Matplotlib}

\cvsection{Awards}

\cvachievement{\faTrophy}{CSIR-NET (JRF)}{Life Sciences, June 2021 --- AIR \textbf{216}}

\cvachievement{\faTrophy}{GATE}{Life Sciences (XL), 2022}

\end{paracol}
\end{document}
